%% network.tex
%%
%% Copyright (C) 2000-2018 Simone Piccardi.  Permission is granted to
%% copy, distribute and/or modify this document under the terms of the GNU Free
%% Documentation License, Version 1.1 or any later version published by the
%% Free Software Foundation; with the Invariant Sections being "Un preambolo",
%% with no Front-Cover Texts, and with no Back-Cover Texts.  A copy of the
%% license is included in the section entitled "GNU Free Documentation
%% License".
%%

\chapter{Introduzione alla programmazione di rete}
\label{cha:network}

In questo capitolo sarà fatta un'introduzione ai concetti generali che servono
come prerequisiti per capire la programmazione di rete, non tratteremo quindi
aspetti specifici ma faremo una breve introduzione ai modelli più comuni usati
nella programmazione di rete, per poi passare ad un esame a grandi linee dei
protocolli di rete e di come questi sono organizzati e interagiscono.

In particolare, avendo assunto l'ottica di un'introduzione mirata alla
programmazione, ci concentreremo sul gruppo di protocolli più diffuso, il
TCP/IP, che è quello che sta alla base di Internet, avendo cura di
sottolineare i concetti più importanti da conoscere per la scrittura dei
programmi.



\section{Modelli di programmazione}
\label{sec:net_prog_model}

La differenza principale fra un'applicazione di rete e un programma normale è
che quest'ultima per definizione concerne la comunicazione fra processi
diversi, che in generale non girano neanche sulla stessa macchina. Questo già
prefigura un cambiamento completo rispetto all'ottica del programma monolitico
all'interno del quale vengono eseguite tutte le istruzioni, e chiaramente
presuppone un sistema operativo multitasking in grado di eseguire più processi
contemporaneamente.

In questa prima sezione esamineremo brevemente i principali modelli di
programmazione in uso. Ne daremo una descrizione assolutamente generica e
superficiale, che ne illustri le caratteristiche principali, non essendo fra
gli scopi del testo approfondire questi argomenti.

\subsection{Il modello \textit{client-server}}
\label{sec:net_cliserv}

L'architettura fondamentale su cui si basa gran parte della programmazione di
rete sotto Linux (e sotto Unix in generale) è il modello
\textit{client-server} caratterizzato dalla presenza di due categorie di
soggetti, i programmi di servizio, chiamati \textit{server}, che ricevono le
richieste e forniscono le risposte, ed i programmi di utilizzo, detti
\textit{client}.

In generale un server può (di norma deve) essere in grado di rispondere a più
di un client, per cui è possibile che molti programmi possano interagire
contemporaneamente, quello che contraddistingue il modello però è che
l'architettura dell'interazione è sempre nei termini di molti verso uno, il
server, che viene ad assumere un ruolo privilegiato.

Seguono questo modello tutti i servizi fondamentali di Internet, come le
pagine web, la posta elettronica, ftp, telnet, ssh e praticamente ogni
servizio che viene fornito tramite la rete, anche se, come abbiamo visto, il
modello è utilizzato in generale anche per programmi che non fanno
necessariamente uso della rete, come gli esempi che abbiamo usato in
cap.~\ref{cha:IPC} a proposito della comunicazione fra processi nello stesso
sistema.

Normalmente si dividono i server in due categorie principali, e vengono detti
\textsl{concorrenti} o \textsl{iterativi}, sulla base del loro comportamento.
Un \textsl{server iterativo} risponde alla richiesta inviando i dati e resta
occupato e non rispondendo ad ulteriori richieste fintanto che non ha fornito
una risposta alla richiesta. Una volta completata la risposta il server
diventa di nuovo disponibile.

Un \textsl{server concorrente} al momento di trattare la richiesta crea un
processo figlio (o un \textit{thread}) incaricato di fornire i servizi
richiesti, per porsi immediatamente in attesa di ulteriori richieste. In
questo modo, con sistemi multitasking, più richieste possono essere
soddisfatte contemporaneamente. Una volta che il processo figlio ha concluso
il suo lavoro esso di norma viene terminato, mentre il server originale resta
sempre attivo.


\subsection{Il modello \textit{peer-to-peer}}
\label{sec:net_peertopeer}

Come abbiamo visto il tratto saliente dell'architettura \textit{client-server}
è quello della preminenza del server rispetto ai client, le architetture
\textit{peer-to-peer} si basano su un approccio completamente opposto che è
quello di non avere nessun programma che svolga un ruolo preminente.

Questo vuol dire che in generale ciascun programma viene ad agire come un nodo
in una rete potenzialmente paritetica; ciascun programma si trova pertanto a
ricevere ed inviare richieste ed a ricevere ed inviare risposte, e non c'è più
la separazione netta dei compiti che si ritrova nelle architetture
\textit{client-server}.

Le architetture \textit{peer-to-peer} sono salite alla ribalta con
l'esplosione del fenomeno Napster, ma gli stessi protocolli di routing sono un
buon esempio di architetture \textit{peer-to-peer}, in cui ciascun nodo,
tramite il demone che gestisce il routing, richiede ed invia informazioni ad
altri nodi.

In realtà in molti casi di architetture classificate come
\textit{peer-to-peer} non è detto che la struttura sia totalmente paritetica e
ci sono parecchi esempi in cui alcuni servizi vengono centralizzati o
distribuiti gerarchicamente, come avveniva per lo stesso Napster, in cui le
ricerche erano effettuate su un server centrale.



\subsection{Il modello \textit{three-tier}}
\label{sec:net_three_tier}

Benché qui sia trattato a parte, il modello \textit{three-tier} in realtà è
una estensione del modello \textit{client-server}. Con il crescere della
quantità dei servizi forniti in rete (in particolare su Internet) ed al numero
di accessi richiesto. Si è così assistito anche ad una notevole crescita di
complessità, in cui diversi servizi venivano ad essere integrati fra di loro.

In particolare sempre più spesso si assiste ad una integrazione di servizi di
database con servizi di web, in cui le pagine vengono costruite dinamicamente
sulla base dei dati contenuti nel database. In tutti questi casi il problema
fondamentale di una architettura \textit{client-server} è che la richiesta di
un servizio da parte di un gran numero di client si scontra con il collo di
bottiglia dell'accesso diretto ad un unico server, con gravi problemi di
scalabilità.

Rispondere a queste esigenze di scalabilità il modello più semplice (chiamato
talvolta \textit{two-tier}) da adottare è stata quello di distribuire il
carico delle richieste su più server identici, mantenendo quindi
sostanzialmente inalterata l'architettura \textit{client-server} originale.

Nel far questo ci si scontra però con gravi problemi di manutenibilità dei
servizi, in particolare per quanto riguarda la sincronizzazione dei dati, e di
inefficienza dell'uso delle risorse. Il problema è particolarmente grave ad
esempio per i database che non possono essere replicati e sincronizzati
facilmente, e che sono molto onerosi, la loro replicazione è costosa e
complessa.

È a partire da queste problematiche che nasce il modello \textit{three-tier},
che si struttura, come dice il nome, su tre livelli. Il primo livello, quello
dei client che eseguono le richieste e gestiscono l'interfaccia con l'utente,
resta sostanzialmente lo stesso del modello \textit{client-server}, ma la
parte server viene suddivisa in due livelli, introducendo un
\textit{middle-tier}, su cui deve appoggiarsi tutta la logica di analisi delle
richieste dei client per ottimizzare l'accesso al terzo livello, che è quello
che si limita a fornire i dati dinamici che verranno usati dalla logica
implementata nel \textit{middle-tier} per eseguire le operazioni richieste dai
client.

In questo modo si può disaccoppiare la logica dai dati, replicando la prima,
che è molto meno soggetta a cambiamenti ed evoluzione, e non soffre di
problemi di sincronizzazione, e centralizzando opportunamente i secondi. In
questo modo si può distribuire il carico ed accedere in maniera efficiente i
dati.


\subsection{Il modello \textit{broadcast}}
\label{sec:net_broadcast}

Uno specifico modello relativo alla programmazione di rete è poi quello in cui
è possibile, invece della classica comunicazione uno ad uno comunque usata in
tutti i modelli precedenti (anche nel \textit{peer-to-peer} la comunicazione è
comunque fra singoli ``\textit{peer}''), una comunicazione da uno a molti.

\itindbeg{broadcast}

Questo modello nasce dal fatto che molte tecnologie di rete (ed in particolare
Ethernet, che è probabilmente la più diffusa) hanno il supporto per effettuare
una comunicazione in cui un nodo qualunque della rete più inviare informazioni
in contemporanea a tutti gli altri. In questo caso si parla di
\textit{broadcast}, utilizzando la nomenclatura usata per le trasmissioni
radio, anche se in realtà questo tipo di comunicazione è eseguibile da un nodo
qualunque per cui tutti quanti possono ricoprire sia il ruolo di trasmettitore
che quello di ricevitore.

\itindbeg{multicast}

In genere si parla di \textit{broadcast} quando la trasmissione uno a molti è
possibile fra qualunque nodo di una rete e gli altri, ed è supportata
direttamente dalla tecnologia di collegamento utilizzata. L'utilizzo di questa
forma di comunicazione da uno a molti però può risultare molto utile anche
quando questo tipo di supporto non è disponibile (come ad esempio su Internet,
dove non si possono contattare tutti i nodi presenti). 

\itindend{broadcast}

In tal caso alcuni protocolli di rete (e quelli usati per Internet sono fra
questi) supportano una variante del\textit{broadcast}, detta
\textit{multicast}, in cui resta possibile fare una comunicazione uno a molti,
in cui una applicazione invia i pacchetti a molte altre, in genere passando
attraverso un opportuno supporto degli apparati ed una qualche forma di
registrazione che consente la distribuzione della cominicazione ai nodi
interessati. 

\itindend{multicast}

Ovviamente i programmi che devono realizzare un tipo di comunicazione di
questo tipo (come ad esempio potrebbero essere quelli che effettuano uno
\textit{streaming} di informazioni) devono rispondere a delle problematiche
del tutto diverse da quelle classiche illustrate nei modelli precedenti, e
costituiscono pertanto un'altra classe completamente a parte.


\section{I protocolli di rete}
\label{sec:net_protocols}

Parlando di reti di computer si parla in genere di un insieme molto vasto ed
eterogeneo di mezzi di comunicazione che vanno dal cavo telefonico, alla fibra
ottica, alle comunicazioni via satellite o via radio; per rendere possibile la
comunicazione attraverso un così variegato insieme di mezzi sono stati
adottati molti protocolli, il più famoso dei quali, quello alla base del
funzionamento di Internet, è il gruppo di protocolli comunemente chiamato
TCP/IP.

\subsection{Il modello ISO/OSI}
\label{sec:net_iso_osi}

Una caratteristica comune dei protocolli di rete è il loro essere strutturati
in livelli sovrapposti; in questo modo ogni protocollo di un certo livello
realizza le sue funzionalità basandosi su un protocollo del livello
sottostante.  Questo modello di funzionamento è stato standardizzato dalla
\textit{International Standards Organization} (ISO) che ha preparato fin dal
1984 il Modello di Riferimento \textit{Open Systems Interconnection} (OSI),
strutturato in sette livelli, secondo quanto riportato in
tab.~\ref{tab:net_osilayers}.

\begin{table}[htb]
  \centering
  \begin{tabular}{|l|c|c|} 
    \hline
    \textbf{Livello} & \multicolumn{2}{|c|}{\textbf{Nome}} \\
    \hline
    \hline
    Livello 7&\textit{Application}  &\textsl{Applicazione}\\ 
    Livello 6&\textit{Presentation} &\textsl{Presentazione} \\ 
    Livello 5&\textit{Session}      &\textsl{Sessione} \\ 
    Livello 4&\textit{Transport}    &\textsl{Trasporto} \\ 
    Livello 3&\textit{Network}      &\textsl{Rete}\\ 
    Livello 2&\textit{DataLink}     &\textsl{Collegamento Dati} \\
    Livello 1&\textit{Physical}   &\textsl{Connessione Fisica} \\
    \hline
\end{tabular}
\caption{I sette livelli del protocollo ISO/OSI.}
\label{tab:net_osilayers}
\end{table}

Il modello ISO/OSI è stato sviluppato in corrispondenza alla definizione della
serie di protocolli X.25 per la commutazione di pacchetto; come si vede è un
modello abbastanza complesso\footnote{infatti per memorizzarne i vari livelli
  è stata creata la frase \textit{All people seem to need data processing}, in
  cui ciascuna parola corrisponde all'iniziale di uno dei livelli.}, tanto che
usualmente si tende a suddividerlo in due parti, secondo lo schema mostrato in
fig.~\ref{fig:net_osi_tcpip_comp}, con un \textit{upper layer} che riguarda
solo le applicazioni, che viene realizzato in \textit{user space}, ed un
\textit{lower layer} in cui si mescolano la gestione fatta dal kernel e le
funzionalità fornite dall'hardware.

Il modello ISO/OSI mira ad effettuare una classificazione completamente
generale di ogni tipo di protocollo di rete; nel frattempo però era stato
sviluppato anche un altro modello, relativo al protocollo TCP/IP, che è quello
su cui è basata Internet, che è diventato uno standard de facto.  Questo
modello viene talvolta chiamato anche modello \textit{DoD} (sigla che sta per
\textit{Department of Defense}), dato che fu sviluppato dall'agenzia ARPA per
il Dipartimento della Difesa Americano.

\begin{figure}[!htb]
  \centering
  \includegraphics[width=12cm]{img/iso_tcp_comp}
  \caption{Struttura a livelli dei protocolli OSI e TCP/IP, con la relative
    corrispondenze e la divisione fra \textit{kernel space} e \textit{user
      space}.}
  \label{fig:net_osi_tcpip_comp}
\end{figure}

La scelta fra quale dei due modelli utilizzare dipende per lo più dai gusti
personali. Come caratteristiche generali il modello ISO/OSI è più teorico e
generico, basato separazioni funzionali, mentre il modello TCP/IP è più vicino
alla separazione concreta dei vari strati del sistema operativo; useremo
pertanto quest'ultimo, anche per la sua maggiore semplicità. Questa semplicità
ha un costo quando si fa riferimento agli strati più bassi, che sono in
effetti descritti meglio dal modello ISO/OSI, in quanto gran parte dei
protocolli di trasmissione hardware sono appunto strutturati sui due livelli
di \textit{Data Link} e \textit{Connection}.


\subsection{Il modello TCP/IP (o DoD)}
\label{sec:net_tcpip_overview}

Così come ISO/OSI anche il modello del TCP/IP è stato strutturato in livelli
(riassunti in tab.~\ref{tab:net_layers}); un confronto fra i due è riportato
in fig.~\ref{fig:net_osi_tcpip_comp} dove viene evidenziata anche la
corrispondenza fra i rispettivi livelli (che comunque è approssimativa) e su
come essi vanno ad inserirsi all'interno del sistema rispetto alla divisione
fra \textit{user space} e \textit{kernel space} spiegata in
sez.~\ref{sec:intro_unix_struct}.\footnote{in realtà è sempre possibile
  accedere dallo \textit{user space}, attraverso una opportuna interfaccia
  (come vedremo in sez.~\ref{sec:sock_sa_packet}), ai livelli inferiori del
  protocollo.}

\begin{table}[htb]
  \centering
  \begin{tabular}{|l|c|c|l|} 
    \hline
    \textbf{Livello} & \multicolumn{2}{|c|}{\textbf{Nome}} & \textbf{Esempi} \\
    \hline
    \hline
    Livello 4 & \textit{Application} & \textsl{Applicazione}& 
                                       Telnet, FTP, ecc. \\ 
    Livello 3 & \textit{Transport}   & \textsl{Trasporto} & TCP, UDP\\ 
    Livello 2 & \textit{Network}     & \textsl{Rete}      & IP, (ICMP, IGMP)\\ 
    Livello 1 & \textit{Link}        & \textsl{Collegamento}& 
                                       Device driver \& scheda di interfaccia\\
    \hline
\end{tabular}
\caption{I quattro livelli del protocollo TCP/IP.}
\label{tab:net_layers}
\end{table}

Come si può notare come il modello TCP/IP è più semplice del modello ISO/OSI
ed è strutturato in soli quattro livelli. Il suo nome deriva dai due
principali protocolli che lo compongono, il TCP (\textit{Trasmission Control
  Protocol}) che copre il livello 3 e l'IP (\textit{Internet Protocol}) che
copre il livello 2. Le funzioni dei vari livelli sono le seguenti:

\begin{basedescript}{\desclabelwidth{2.5cm}\desclabelstyle{\nextlinelabel}}
\item[\textbf{Applicazione}] É relativo ai programmi di interfaccia con la
  rete, in genere questi vengono realizzati secondo il modello client-server
  (vedi sez.~\ref{sec:net_cliserv}), realizzando una comunicazione secondo un
  protocollo che è specifico di ciascuna applicazione.
\item[\textbf{Trasporto}] Fornisce la comunicazione tra le due stazioni
  terminali su cui girano gli applicativi, regola il flusso delle
  informazioni, può fornire un trasporto affidabile, cioè con recupero degli
  errori o inaffidabile. I protocolli principali di questo livello sono il TCP
  e l'UDP.
\item[\textbf{Rete}] Si occupa dello smistamento dei singoli pacchetti su una
  rete complessa e interconnessa, a questo stesso livello operano i protocolli
  per il reperimento delle informazioni necessarie allo smistamento, per lo
  scambio di messaggi di controllo e per il monitoraggio della rete. Il
  protocollo su cui si basa questo livello è IP (sia nella attuale versione,
  IPv4, che nella nuova versione, IPv6).
\item[\textbf{Collegamento}] È responsabile per l'interfacciamento al
  dispositivo elettronico che effettua la comunicazione fisica, gestendo
  l'invio e la ricezione dei pacchetti da e verso l'hardware.
\end{basedescript}

La comunicazione fra due stazioni remote avviene secondo le modalità
illustrate in fig.~\ref{fig:net_tcpip_data_flux}, dove si è riportato il flusso
dei dati reali e i protocolli usati per lo scambio di informazione su ciascun
livello. Si è genericamente indicato \textit{ethernet} per il livello 1, anche
se in realtà i protocolli di trasmissione usati possono essere molti altri.

\begin{figure}[!htb]
  \centering \includegraphics[width=13cm]{img/tcp_data_flux}
  \caption{Strutturazione del flusso dei dati nella comunicazione fra due
    applicazioni attraverso i protocolli della suite TCP/IP.}
  \label{fig:net_tcpip_data_flux}
\end{figure}

Per chiarire meglio la struttura della comunicazione attraverso i vari
protocolli mostrata in fig.~\ref{fig:net_tcpip_data_flux}, conviene prendere in
esame i singoli passaggi fatti per passare da un livello al sottostante,
la procedura si può riassumere nei seguenti passi:
\begin{itemize}
\item Le singole applicazioni comunicano scambiandosi i dati ciascuna secondo
  un suo specifico formato. Per applicazioni generiche, come la posta o le
  pagine web, viene di solito definito ed implementato quello che viene
  chiamato un protocollo di applicazione (esempi possono essere HTTP, POP,
  SMTP, ecc.), ciascuno dei quali è descritto in un opportuno standard, di
  solito attraverso un RFC (l'acronimo RFC sta per
  \itindex{Request~For~Comment~(RFC)} \textit{Request For Comment} ed è la
  procedura attraverso la quale vengono proposti gli standard per Internet).
\item I dati delle applicazioni vengono inviati al livello di trasporto usando
  un'interfaccia opportuna (i \textit{socket}, che esamineremo in dettaglio in
  cap.~\ref{cha:socket_intro}). Qui verranno spezzati in pacchetti di
  dimensione opportuna e inseriti nel protocollo di trasporto, aggiungendo ad
  ogni pacchetto le informazioni necessarie per la sua gestione. Questo
  processo viene svolto direttamente nel kernel, ad esempio dallo stack TCP,
  nel caso il protocollo di trasporto usato sia questo.
\item Una volta composto il pacchetto nel formato adatto al protocollo di
  trasporto usato questo sarà passato al successivo livello, quello di rete,
  che si occupa di inserire le opportune informazioni per poter effettuare
  l'instradamento nella rete ed il recapito alla destinazione finale. In
  genere questo è il livello di IP (Internet Protocol), a cui vengono inseriti
  i numeri IP che identificano i computer su Internet.
\item L'ultimo passo è il trasferimento del pacchetto al driver della
  interfaccia di trasmissione, che si incarica di incapsularlo nel relativo
  protocollo di trasmissione. Questo può avvenire sia in maniera diretta, come
  nel caso di ethernet, in cui i pacchetti vengono inviati sulla linea
  attraverso le schede di rete, che in maniera indiretta con protocolli come
  PPP o SLIP, che vengono usati come interfaccia per far passare i dati su
  altri dispositivi di comunicazione (come la seriale o la parallela).
\end{itemize}


\subsection{Criteri generali dell'architettura del TCP/IP}
\label{sec:net_tcpip_design}

La filosofia architetturale del TCP/IP è semplice: costruire una rete che
possa sopportare il carico in transito, ma permettere ai singoli nodi di
scartare pacchetti se il carico è temporaneamente eccessivo, o se risultano
errati o non recapitabili.

L'incarico di rendere il recapito pacchetti affidabile non spetta al livello
di rete, ma ai livelli superiori. Pertanto il protocollo IP è per sua natura
inaffidabile, in quanto non è assicurata né una percentuale di successo né un
limite sui tempi di consegna dei pacchetti.

È il livello di trasporto che si deve occupare (qualora necessiti) del
controllo del flusso dei dati e del recupero degli errori; questo è realizzato
dal protocollo TCP. La sede principale di "\textit{intelligenza}" della rete è
pertanto al livello di trasporto o ai livelli superiori.

Infine le singole stazioni collegate alla rete non fungono soltanto da punti
terminali di comunicazione, ma possono anche assumere il ruolo di
\textit{router} (\textsl{instradatori}), per l'interscambio di pacchetti da
una rete ad un'altra. Questo rende possibile la flessibilità della rete che è
in grado di adattarsi ai mutamenti delle interconnessioni.

La caratteristica essenziale che rende tutto ciò possibile è la strutturazione
a livelli tramite l'incapsulamento. Ogni pacchetto di dati viene incapsulato
nel formato del livello successivo, fino al livello del collegamento fisico.
In questo modo il pacchetto ricevuto ad un livello \textit{n} dalla stazione
di destinazione è esattamente lo stesso spedito dal livello \textit{n} dalla
sorgente.  Questo rende facile il progettare il software facendo riferimento
unicamente a quanto necessario ad un singolo livello, con la confidenza che
questo poi sarà trattato uniformemente da tutti i nodi della rete.


\section{La struttura del TCP/IP}
\label{sec:net_tpcip}

Come accennato in sez.~\ref{sec:net_protocols} il TCP/IP è un insieme di
protocolli diversi, che operano su 4 livelli diversi. Per gli interessi della
programmazione di rete però sono importanti principalmente i due livelli
centrali, e soprattutto quello di trasporto.

La principale interfaccia usata nella programmazione di rete, quella dei
socket (che vedremo in sez.~\ref{cha:socket_intro}), è infatti un'interfaccia
nei confronti di quest'ultimo.  Questo avviene perché al di sopra del livello
di trasporto i programmi hanno a che fare solo con dettagli specifici delle
applicazioni, mentre al di sotto vengono curati tutti i dettagli relativi alla
comunicazione. È pertanto naturale definire una interfaccia di programmazione
su questo confine, tanto più che è proprio lì (come evidenziato in
fig.~\ref{fig:net_osi_tcpip_comp}) che nei sistemi Unix (e non solo) viene
inserita la divisione fra \textit{kernel space} e \textit{user space}.

In realtà in un sistema Unix è possibile accedere anche agli altri livelli (e
non solo a quello di trasporto) con opportune interfacce di programmazione
(vedi sez.~\ref{sec:sock_sa_packet}), ma queste vengono usate solo quando si
debbano fare applicazioni di sistema per il controllo della rete a basso
livello, di uso quindi molto specialistico.

In questa sezione daremo una descrizione sommaria dei vari protocolli del
TCP/IP, concentrandoci, per le ragioni appena esposte, sul livello di
trasporto.  All'interno di quest'ultimo privilegeremo poi il protocollo TCP,
per il ruolo centrale che svolge nella maggior parte delle applicazioni.


\subsection{Il quadro generale}
\label{sec:net_tcpip_general}

Benché si parli di TCP/IP questa famiglia di protocolli è composta anche da
molti membri. In fig.~\ref{fig:net_tcpip_overview} si è riportato uno schema
che mostra un panorama sui principali protocolli della famiglia, e delle loro
relazioni reciproche e con alcune dalle principali applicazioni che li usano.

\begin{figure}[!htb]
  \centering
  \includegraphics[width=13cm]{img/tcpip_overview}  
  \caption{Panoramica sui vari protocolli che compongono la suite TCP/IP.}
  \label{fig:net_tcpip_overview}
\end{figure}

I vari protocolli riportati in fig.~\ref{fig:net_tcpip_overview} sono i
seguenti:
\begin{basedescript}{\desclabelwidth{1.7cm}\desclabelstyle{\nextlinelabel}}
\item[\textsl{IPv4}] \textit{Internet Protocol version 4}. È quello che
  comunemente si chiama IP. Ha origine negli anni '80 e da allora è la base su
  cui è costruita Internet. Usa indirizzi a 32 bit, e mantiene tutte le
  informazioni di instradamento e controllo per la trasmissione dei pacchetti
  sulla rete; tutti gli altri protocolli della suite (eccetto ARP e RARP, e
  quelli specifici di IPv6) vengono trasmessi attraverso di esso.
\item[\textsl{IPv6}] \textit{Internet Protocol version 6}. È stato progettato
  a metà degli anni '90 per rimpiazzare IPv4. Ha uno spazio di indirizzi
  ampliato 128 bit che consente più gerarchie di indirizzi,
  l'auto-configurazione, ed un nuovo tipo di indirizzi, gli \textit{anycast},
  che consentono di inviare un pacchetto ad una stazione su un certo gruppo.
  Effettua lo stesso servizio di trasmissione dei pacchetti di IPv4 di cui
  vuole essere un sostituto.
\item[\textsl{TCP}] \textit{Trasmission Control Protocol}. È un protocollo
  orientato alla connessione che provvede un trasporto affidabile per un
  flusso di dati bidirezionale fra due stazioni remote. Il protocollo ha cura
  di tutti gli aspetti del trasporto dei dati, come l'\textit{acknowledgment}
  (il ricevuto), i timeout, la ritrasmissione, ecc. È usato dalla maggior
  parte delle applicazioni.
\item[\textsl{UDP}] \textit{User Datagram Protocol}. È un protocollo senza
  connessione, per l'invio di dati a pacchetti. Contrariamente al TCP il
  protocollo non è affidabile e non c'è garanzia che i pacchetti raggiungano
  la loro destinazione, si perdano, vengano duplicati, o abbiano un
  particolare ordine di arrivo.
\item[\textsl{ICMP}] \textit{Internet Control Message Protocol}. È il
  protocollo usato a livello 2 per gestire gli errori e trasportare le
  informazioni di controllo fra stazioni remote e instradatori (cioè fra
  \textit{host} e \textit{router}). I messaggi sono normalmente generati dal
  software del kernel che gestisce la comunicazione TCP/IP, anche se ICMP può
  venire usato direttamente da alcuni programmi come \cmd{ping}. A volte ci
  si riferisce ad esso come ICPMv4 per distinguerlo da ICMPv6.
\item[\textsl{IGMP}] \textit{Internet Group Management Protocol}. É un
  protocollo di livello 2 usato per il \textit{multicast} (vedi
  sez.~\ref{sec:xxx_multicast}).  Permette alle stazioni remote di notificare
  ai router che supportano questa comunicazione a quale gruppo esse
  appartengono.  Come ICMP viene implementato direttamente sopra IP.
\item[\textsl{ARP}] \textit{Address Resolution Protocol}. È il protocollo che
  mappa un indirizzo IP in un indirizzo hardware sulla rete locale. È usato in
  reti di tipo \textit{broadcast} come Ethernet, Token Ring o FDDI che hanno
  associato un indirizzo fisico (il \textit{MAC address}) alla interfaccia, ma
  non serve in connessioni punto-punto.
\item[\textsl{RARP}] \textit{Reverse Address Resolution Protocol}. È il
  protocollo che esegue l'operazione inversa rispetto ad ARP (da cui il nome)
  mappando un indirizzo hardware in un indirizzo IP. Viene usato a volte per
  durante l'avvio per assegnare un indirizzo IP ad una macchina.
\item[\textsl{ICMPv6}] \textit{Internet Control Message Protocol, version 6}.
  Combina per IPv6 le funzionalità di ICMPv4, IGMP e ARP.
\item[\textsl{EGP}] \textit{Exterior Gateway Protocol}. È un protocollo di
  routing usato per comunicare lo stato fra gateway vicini a livello di
  \textsl{sistemi autonomi} (vengono chiamati \textit{autonomous
      systems} i raggruppamenti al livello più alto della rete), con
  meccanismi che permettono di identificare i vicini, controllarne la
  raggiungibilità e scambiare informazioni sullo stato della rete. Viene
  implementato direttamente sopra IP. 
\item[\textsl{OSPF}] \textit{Open Shortest Path First}. È in protocollo di
  routing per router su reti interne, che permette a questi ultimi di
  scambiarsi informazioni sullo stato delle connessioni e dei legami che
  ciascuno ha con gli altri. Viene implementato direttamente sopra IP.
\item[\textsl{GRE}] \textit{Generic Routing Encapsulation}. È un protocollo
  generico di incapsulamento che permette di incapsulare un qualunque altro
  protocollo all'interno di IP. 
\item[\textsl{AH}] \textit{Authentication Header}. Provvede l'autenticazione
  dell'integrità e dell'origine di un pacchetto. È una opzione nativa in IPv6
  e viene implementato come protocollo a sé su IPv4. Fa parte della suite di
  IPSEC che provvede la trasmissione cifrata ed autenticata a livello IP.
\item[\textsl{ESP}] \textit{Encapsulating Security Payload}. Provvede la
  cifratura insieme all'autenticazione dell'integrità e dell'origine di un
  pacchetto. Come per AH è opzione nativa in IPv6 e viene implementato come
  protocollo a sé su IPv4.
\item[\textsl{PPP}] \textit{Point-to-Point Protocol}. È un protocollo a
  livello 1 progettato per lo scambio di pacchetti su connessioni punto punto.
  Viene usato per configurare i collegamenti, definire i protocolli di rete
  usati ed incapsulare i pacchetti di dati. È un protocollo complesso con
  varie componenti.
\item[\textsl{SLIP}] \textit{Serial Line over IP}. È un protocollo di livello
  1 che permette di trasmettere un pacchetto IP attraverso una linea seriale.
\end{basedescript}

Gran parte delle applicazioni comunicano usando TCP o UDP, solo alcune, e per
scopi particolari si rifanno direttamente ad IP (ed i suoi correlati ICMP e
IGMP); benché sia TCP che UDP siano basati su IP e sia possibile intervenire a
questo livello con i \textit{raw socket} questa tecnica è molto meno diffusa e
a parte applicazioni particolari si preferisce sempre usare i servizi messi a
disposizione dai due protocolli precedenti.  Per questo, motivo a parte alcuni
brevi accenni su IP in questa sezione, ci concentreremo sul livello di
trasporto.

\subsection{Internet Protocol (IP)}
\label{sec:net_ip}

Quando si parla di \textit{Internet Protocol} (IP) si fa in genere riferimento
ad una versione (la quarta, da cui il nome IPv4) che è quella più usata
comunemente, anche se ormai si sta diffondendo sempre di più la nuova versione
IPv6. Il protocollo IPv4 venne standardizzato nel 1981
dall'\href{http://www.ietf.org/rfc/rfc0719.txt}{RFC~719}.

Il protocollo IP (indipendentemente dalla versione) nasce per disaccoppiare le
applicazioni della struttura hardware delle reti di trasmissione, e creare una
interfaccia di trasmissione dei dati indipendente dal sottostante substrato di
interconnessione fisica, che può essere realizzato con le tecnologie più
disparate (Ethernet, Token Ring, FDDI, ecc.).  Il compito di IP è pertanto
quello di trasmettere i pacchetti da un computer all'altro della rete; le
caratteristiche essenziali con cui questo viene realizzato in sono due:

\begin{itemize}
\item \textit{Universal addressing} la comunicazione avviene fra due stazioni
  remote identificate univocamente con un indirizzo a 32 bit che può
  appartenere ad una sola interfaccia di rete.
\item \textit{Best effort} viene assicurato il massimo impegno nella
  trasmissione, ma non c'è nessuna garanzia per i livelli superiori né sulla
  percentuale di successo né sul tempo di consegna dei pacchetti di dati.
\end{itemize}

Negli anni '90 la crescita vertiginosa del numero di macchine connesse a
Internet ha iniziato a far emergere i vari limiti di IPv4, per risolverne i
problemi si è perciò definita una nuova versione del protocollo, che (saltando
un numero) è diventata la versione 6. IPv6 nasce quindi come evoluzione di
IPv4, mantenendone inalterate le funzioni che si sono dimostrate valide,
eliminando quelle inutili e aggiungendone poche altre per mantenere il
protocollo il più snello e veloce possibile.

I cambiamenti apportati sono comunque notevoli e si possono essere riassunti a
grandi linee nei seguenti punti:
\begin{itemize}
\item l'espansione delle capacità di indirizzamento e instradamento, per
  supportare una gerarchia con più livelli di indirizzamento, un numero di
  nodi indirizzabili molto maggiore e una auto-configurazione degli indirizzi.
\item l'introduzione un nuovo tipo di indirizzamento, l'\textit{anycast} che
  si aggiunge agli usuali \textit{unicast} e \textit{multicast}.
\item la semplificazione del formato dell'intestazione (\textit{header}) dei
  pacchetti, eliminando o rendendo opzionali alcuni dei campi di IPv4, per
  eliminare la necessità di rielaborazione della stessa da parte dei router e
  contenere l'aumento di dimensione dovuto all'ampliamento degli indirizzi.
\item un supporto per le opzioni migliorato, per garantire una trasmissione
  più efficiente del traffico normale, limiti meno stringenti sulle dimensioni
  delle opzioni, e la flessibilità necessaria per introdurne di nuove in
  futuro.
\item il supporto per delle capacità di \textsl{qualità di servizio} (QoS) che
  permettano di identificare gruppi di dati per i quali si può provvedere un
  trattamento speciale (in vista dell'uso di Internet per applicazioni
  multimediali e/o ``real-time'').
\end{itemize}

Maggiori dettagli riguardo a caratteristiche, notazioni e funzionamento del
protocollo IP sono forniti nell'appendice sez.~\ref{sec:ip_protocol}.

 
\subsection{User Datagram Protocol (UDP)}
\label{sec:net_udp}

Il protocollo UDP è un protocollo di trasporto molto semplice; la sua
descrizione completa è contenuta
dell'\href{http://www.ietf.org/rfc/rfc0768.txt}{RFC~768}, ma in sostanza esso
è una semplice interfaccia al protocollo IP dal livello di trasporto. Quando
un'applicazione usa UDP essa scrive un pacchetto di dati (il cosiddetto
\textit{datagram} che da il nome al protocollo) su un socket, al pacchetto
viene aggiunto un header molto semplice (per una descrizione più accurata vedi
sez.~\ref{sec:udp_protocol}), e poi viene passato al livello superiore (IPv4 o
IPv6 che sia) che lo spedisce verso la destinazione.  Dato che né IPv4 né IPv6
garantiscono l'affidabilità niente assicura che il pacchetto arrivi a
destinazione, né che più pacchetti arrivino nello stesso ordine in cui sono
stati spediti.

Pertanto il problema principale che si affronta quando si usa UDP è la
mancanza di affidabilità, se si vuole essere sicuri che i pacchetti arrivino a
destinazione occorrerà provvedere con l'applicazione, all'interno della quale
si dovrà inserire tutto quanto necessario a gestire la notifica di
ricevimento, la ritrasmissione, il timeout. 

Si tenga conto poi che in UDP niente garantisce che i pacchetti arrivino nello
stesso ordine in cui sono stati trasmessi, e può anche accadere che i
pacchetti vengano duplicati nella trasmissione, e non solo perduti. Di tutto
questo di nuovo deve tenere conto l'applicazione.

Un altro aspetto di UDP è che se un pacchetto raggiunge correttamente la
destinazione esso viene passato all'applicazione ricevente in tutta la sua
lunghezza, la trasmissione avviene perciò per \textit{record} la cui lunghezza
viene anche essa trasmessa all'applicazione all'atto del ricevimento.

Infine UDP è un protocollo che opera senza connessione
(\textit{connectionless}) in quanto non è necessario stabilire nessun tipo di
relazione tra origine e destinazione dei pacchetti. Si hanno così situazioni
in cui un client può scrivere su uno stesso socket pacchetti destinati a
server diversi, o un server ricevere su un socket pacchetti provenienti da
client diversi.  Il modo più semplice di immaginarsi il funzionamento di UDP è
quello della radio, in cui si può \textsl{trasmettere} e \textsl{ricevere} da
più stazioni usando la stessa frequenza.

Nonostante gli evidenti svantaggi comportati dall'inaffidabilità UDP ha il
grande pregio della velocità, che in certi casi è essenziale; inoltre si
presta bene per le applicazioni in cui la connessione non è necessaria, e
costituirebbe solo un peso in termini di prestazioni, mentre una perdita di
pacchetti può essere tollerata: ad esempio le applicazioni di streaming e
quelle che usano il \textit{multicast}.

\subsection{Transport Control Protocol (TCP)}
\label{sec:net_tcp}

Il TCP è un protocollo molto complesso, definito
nell'\href{http://www.ietf.org/rfc/rfc0739.txt}{RFC~739} e completamente
diverso da UDP; alla base della sua progettazione infatti non stanno
semplicità e velocità, ma la ricerca della massima affidabilità possibile
nella trasmissione dei dati.

La prima differenza con UDP è che TCP provvede sempre una connessione diretta
fra un client e un server, attraverso la quale essi possono comunicare; per
questo il paragone più appropriato per questo protocollo è quello del
collegamento telefonico, in quanto prima viene stabilita una connessione fra
due i due capi della comunicazione su cui poi effettuare quest'ultima.

Caratteristica fondamentale di TCP è l'affidabilità; quando i dati vengono
inviati attraverso una connessione ne viene richiesto un ``\textsl{ricevuto}''
(il cosiddetto \textit{acknowlegment}), se questo non arriva essi verranno
ritrasmessi per un determinato numero di tentativi, intervallati da un periodo
di tempo crescente, fino a che sarà considerata fallita o caduta la
connessione (e sarà generato un errore di \textit{timeout}); il periodo di
tempo dipende dall'implementazione e può variare far i quattro e i dieci
minuti.

Inoltre, per tenere conto delle diverse condizioni in cui può trovarsi la
linea di comunicazione, TCP comprende anche un algoritmo di calcolo dinamico
del tempo di andata e ritorno dei pacchetti fra un client e un server (il
cosiddetto RTT, \textit{Round Trip Time}), che lo rende in grado di adattarsi
alle condizioni della rete per non generare inutili ritrasmissioni o cadere
facilmente in timeout.

Inoltre TCP è in grado di preservare l'ordine dei dati assegnando un numero di
sequenza ad ogni byte che trasmette. Ad esempio se un'applicazione scrive 3000
byte su un socket TCP, questi potranno essere spezzati dal protocollo in due
segmenti (le unità di dati passate da TCP a IP vengono chiamate
\textit{segment}) di 1500 byte, di cui il primo conterrà il numero di sequenza
$1-1500$ e il secondo il numero $1501-3000$. In questo modo anche se i
segmenti arrivano a destinazione in un ordine diverso, o se alcuni arrivano
più volte a causa di ritrasmissioni dovute alla perdita degli
\textit{acknowlegment}, all'arrivo sarà comunque possibile riordinare i dati e
scartare i duplicati.

\itindbeg{advertised~window}

Il protocollo provvede anche un controllo di flusso (\textit{flow control}),
cioè specifica sempre all'altro capo della trasmissione quanti dati può
ricevere tramite una \textit{advertised window} (letteralmente
``\textsl{finestra annunciata}''), che indica lo spazio disponibile nel buffer
di ricezione, cosicché nella trasmissione non vengano inviati più dati di
quelli che possono essere ricevuti.

Questa finestra cambia dinamicamente diminuendo con la ricezione dei dati dal
socket ed aumentando con la lettura di quest'ultimo da parte
dell'applicazione, se diventa nulla il buffer di ricezione è pieno e non
verranno accettati altri dati.  Si noti che UDP non provvede niente di tutto
ciò per cui nulla impedisce che vengano trasmessi pacchetti ad un ritmo che il
ricevente non può sostenere.

\itindend{advertised~window}

Infine attraverso TCP la trasmissione è sempre bidirezionale (in inglese si
dice che è \textit{full-duplex}). È cioè possibile sia trasmettere che
ricevere allo stesso tempo, il che comporta che quanto dicevamo a proposito
del controllo di flusso e della gestione della sequenzialità dei dati viene
effettuato per entrambe le direzioni di comunicazione.

% TODO mettere riferimento alla appendice su TCP quando ci sarà
%% Una descrizione più accurata del protocollo è fornita in appendice
%% sez.~\ref{sec:tcp_protocol}.

\subsection{Limiti e dimensioni riguardanti la trasmissione dei dati}
\label{sec:net_lim_dim}

Un aspetto di cui bisogna tenere conto nella programmazione di rete, e che
ritornerà in seguito, quando tratteremo gli aspetti più avanzati, è che ci sono
una serie di limiti a cui la trasmissione dei dati attraverso i vari livelli
del protocollo deve sottostare; limiti che è opportuno tenere presente perché
in certi casi si possono avere delle conseguenze sul comportamento delle
applicazioni.

Un elenco di questi limiti, insieme ad un breve accenno alle loro origini ed
alle eventuali implicazioni che possono avere, è il seguente:
\begin{itemize}
\item La dimensione massima di un pacchetto IP è di 65535 byte, compresa
  l'intestazione. Questo è dovuto al fatto che la dimensione è indicata da un
  campo apposito nell'header di IP che è lungo 16 bit (vedi
  fig.~\ref{fig:IP_ipv4_head}).
\item La dimensione massima di un pacchetto normale di IPv6 è di 65575 byte;
  il campo apposito nell'header infatti è sempre a 16 bit, ma la dimensione
  dell'header è fissa e di 40 byte e non è compresa nel valore indicato dal
  suddetto campo. Inoltre IPv6 ha la possibilità di estendere la dimensione di
  un pacchetto usando la \textit{jumbo payload option}.
\itindbeg{Maximum~Transfer~Unit~(MTU)}
\item Molte reti fisiche hanno una MTU (\textit{Maximum Transfer Unit}) che
  dipende dal protocollo specifico usato al livello di connessione fisica. Il
  più comune è quello di ethernet che è pari a 1500 byte, una serie di altri
  valori possibili sono riportati in tab.~\ref{tab:net_mtu_values}.
\end{itemize}

Quando un pacchetto IP viene inviato su una interfaccia di rete e le sue
dimensioni eccedono la MTU viene eseguita la cosiddetta
\textit{frammentazione}, i pacchetti cioè vengono suddivisi in blocchi più
piccoli che possono essere trasmessi attraverso l'interfaccia.\footnote{questo
  accade sia per IPv4 che per IPv6, anche se i pacchetti frammentati sono
  gestiti con modalità diverse, IPv4 usa un flag nell'header, IPv6 una
  opportuna opzione, si veda sez.~\ref{sec:ipv6_protocol}.}

\begin{table}[!htb]
  \centering
  \begin{tabular}[c]{|l|c|}
    \hline
    \textbf{Rete} & \textbf{MTU} \\
    \hline
    \hline
    Hyperlink & 65535 \\
    Token Ring IBM (16 Mbit/sec) & 17914 \\
    Token Ring IEEE 802.5 (4 Mbit/sec) & 4464 \\
    FDDI & 4532 \\
    Ethernet & 1500 \\
    X.25 & 576 \\
    \hline
  \end{tabular}
  \caption{Valori della MTU (\textit{Maximum Transfer Unit}) per una serie di
    diverse tecnologie di rete.} 
  \label{tab:net_mtu_values}
\end{table}

%TODO aggiornare la tabella con dati più recenti

\itindbeg{Path~MTU}

La MTU più piccola fra due stazioni viene in genere chiamata \textit{path
  MTU}, che dice qual è la lunghezza massima oltre la quale un pacchetto
inviato da una stazione ad un'altra verrebbe senz'altro frammentato. Si tenga
conto che non è affatto detto che la \textit{path MTU} sia la stessa in
entrambe le direzioni, perché l'instradamento può essere diverso nei due
sensi, con diverse tipologie di rete coinvolte.

Una delle differenze fra IPv4 e IPv6 é che per IPv6 la frammentazione può
essere eseguita solo alla sorgente, questo vuol dire che i router IPv6 non
frammentano i pacchetti che ritrasmettono (anche se possono frammentare i
pacchetti che generano loro stessi), al contrario di quanto fanno i router
IPv4. In ogni caso una volta frammentati i pacchetti possono essere
riassemblati solo alla destinazione.

Nell'header di IPv4 è previsto il flag \texttt{DF} che specifica che il
pacchetto non deve essere frammentato; un router che riceva un pacchetto le
cui dimensioni eccedano quelle dell'MTU della rete di destinazione genererà un
messaggio di errore ICMPv4 di tipo \textit{destination unreachable,
  fragmentation needed but DF bit set}.  Dato che i router IPv6 non possono
effettuare la frammentazione la ricezione di un pacchetto di dimensione
eccessiva per la ritrasmissione genererà sempre un messaggio di errore ICMPv6
di tipo \textit{packet too big}.

Dato che il meccanismo di frammentazione e riassemblaggio dei pacchetti
comporta inefficienza, normalmente viene utilizzato un procedimento, detto
\textit{path MTU discovery} che permette di determinare il \textit{path MTU}
fra due stazioni; per la realizzazione del procedimento si usa il flag
\texttt{DF} di IPv4 e il comportamento normale di IPv6 inviando delle
opportune serie di pacchetti (per i dettagli vedere
l'\href{http://www.ietf.org/rfc/rfc1191.txt}{RFC~1191} per IPv4 e
l'\href{http://www.ietf.org/rfc/rfc1981.txt}{RFC~1981} per IPv6) fintanto che
non si hanno più errori.

Il TCP usa sempre questo meccanismo, che per le implementazioni di IPv4 è
opzionale, mentre diventa obbligatorio per IPv6.  Per IPv6 infatti, non
potendo i router frammentare i pacchetti, è necessario, per poter comunicare,
conoscere da subito il \textit{path MTU}.

\itindend{Path~MTU}

Infine il TCP definisce una \textit{Maximum Segment Size} o MSS (vedi
sez.~\ref{sec:tcp_protocol}) che annuncia all'altro capo della connessione la
dimensione massima del segmento di dati che può essere ricevuto, così da
evitare la frammentazione. Di norma viene impostato alla dimensione della MTU
dell'interfaccia meno la lunghezza delle intestazioni di IP e TCP, in Linux il
default, mantenuto nella costante \constd{TCP\_MSS} è 512.

\itindend{Maximum~Transfer~Unit~(MTU)}


%%% Local Variables: 
%%% mode: latex
%%% TeX-master: "gapil"
%%% End: 

% LocalWords:  TCP multitasking client ftp telnet ssh cap thread peer to three
% LocalWords:  Napster routing tier two middle International Standards Systems
% LocalWords:  Organization Interconnection tab Application Presentation All of
% LocalWords:  Session Transport DataLink Physical people seem need processing
% LocalWords:  fig upper layer lower kernel DoD Department Defense Connection
% LocalWords:  sez UDP ICMP IGMP device Trasmission Control Protocol l'IP l'UDP
% LocalWords:  IPv ethernet SMTP RFC Request For Comment socket stack PPP ARP
% LocalWords:  router instradatori version RARP anycast Di
% LocalWords:  l'acknoweledgment Datagram Message host ping ICPMv ICMPv Group
% LocalWords:  multicast Address Resolution broadcast Token FDDI MAC address DF
% LocalWords:  Reverse EGP Exterior Gateway gateway autonomous systems OSPF GRE
% LocalWords:  Shortest Path First Generic Encapsulation Authentication Header
% LocalWords:  IPSEC ESP Encapsulating Security Payload Point Line over raw QoS
% LocalWords:  dall' Universal addressing Best effort unicast header dell' RTT
% LocalWords:  datagram connectionless streaming nell' acknowlegment trip flow
% LocalWords:  segment control advertised window nell'header dell'header option
% LocalWords:  payload MTU Transfer Unit Hyperlink IBM Mbit sec IEEE path but
% LocalWords:  dell'MTU destination unreachable fragmentation needed packet too
% LocalWords:  big discovery MSS Size
